\subsection{Ausblick}
\label{subsec:ausblick}

Kurzfristig sind zunächst die fehlgeschlagenen aktuellen Tooling-,
Desktopprofil- und Windows-CI-Pfade zu korrigieren und auf einem festen Commit
erfolgreich zu wiederholen. Ergänzend stehen ein vollständiger Compose-Start
mit Health- und Readiness-Nachweis, ein Lauf der Release Validation sowie die
autorisierte Prüfung der Branch Protection aus. Diese Schritte würden die
offene technische Evidenz erweitern, ohne bereits Produktreife zu belegen.

Mittelfristig sind für ein konkretes Produkt Signierung, Notarisierung,
Updater, Credentials, Deployment und Authentifizierung beziehungsweise RBAC
separat auszugestalten und abzunehmen. Als mögliche Weiterentwicklung der
Baseline lassen sich aus den festgestellten Grenzen eine stärkere
Automatisierung gemeinsamer Contracts, eine striktere Capability-Auswahl und
ein kontrolliertes Synchronisationsmodell für bereits generierte Projekte
ableiten.

Langfristig sollte das Produktivitätspotenzial vergleichend über mehrere
Projekte untersucht werden, etwa anhand von Onboarding-Zeit, Zeit bis zum
ersten lauffähigen Build, manuellen Setup-Schritten, Fehlerhäufigkeit und
Pflegeaufwand. Reale Kundenprojekte könnten zugleich prüfen, ob die analytisch
abgeleiteten Profil- und Eignungskriterien tragen und wo zusätzlicher
Anpassungsbedarf entsteht.
